% TOP_PROBLEM: {"id": 8200004, "title": "Bound Entanglement with Negative Partial Transpose", "category_id": 16, "category": {"id": 16, "name": "physics", "display_name": "Mathematical Physics", "description": "Problems at the intersection of mathematics and physics.", "slug": "physics", "order_index": 16, "created_at": "2026-07-31T15:26:25.671Z"}, "status": "open"}
% FIRST_POSTED: 2026-09-25
% ENTRY_KIND: statement_only
% !TeX program = lualatex
% Definition and sources only; this file is not an LLM solution attempt.
\documentclass[11pt]{article}
\usepackage[margin=25mm]{geometry}
\usepackage{fontspec}
\setmainfont{Latin Modern Roman}
\usepackage{amsmath,amssymb,mathtools}
\usepackage{unicode-math}
\setmathfont{Latin Modern Math}
\usepackage{xurl,hyperref}
\hypersetup{colorlinks=true,urlcolor=blue}
\setcounter{secnumdepth}{0}
\setlength{\parindent}{0pt}
\setlength{\parskip}{5pt}
\setlength{\emergencystretch}{3em}
\begin{document}
\section*{\texorpdfstring{Bound Entanglement with Negative Partial Transpose}{Bound Entanglement with Negative Partial Transpose}}\label{8200004}
\subsection{Definitions and mathematical statement}
A finite-dimensional bipartite state is a positive semidefinite trace-one operator $\rho$ on $H_A\otimes H_B$. It is entangled if it is not a convex combination of product states. Partial transpose $\rho^{T_B}$ transposes the second subsystem in an orthonormal basis; having a negative eigenvalue is basis independent. LOCC consists of local quantum operations with classical communication. The target asks for $\rho^{T_B}\not\succeq0$ but zero asymptotic rate of distilling maximally entangled pairs from $\rho^{\otimes n}$. Equivalently, in the standard finite-dimensional distillability criterion, seek an NPT state with
\[
 \langle\psi|(\rho^{T_B})^{\otimes n}|\psi\rangle\ge0
\]
for every $n\ge1$ and every vector $\psi$ of Schmidt rank at most two across $A^n:B^n$. This is not merely failure of single-copy distillation.
\par\smallskip\textit{Formulation and concept references:} \href{https://arxiv.org/pdf/2002.03233}{[S1]}.

\subsection{Short English statement}
Can a quantum state have a negative partial transpose yet yield no distillable entanglement, even when arbitrarily many copies and two-way classical communication are available?
\subsection{Sources}
\begin{itemize}
\item[S1]P. Horodecki, L. Rudnicki, and K. Zyczkowski, Five Open Problems in Quantum Information Theory, PRX Quantum 3 (2022).
\url{https://arxiv.org/pdf/2002.03233}.
\textit{Formulation source.}
\end{itemize}
\end{document}
